\chapter{Входной гейт рангоизменяющей суперсвязности}

\section{Цель проверки}

Финал Тома VI разрешает новую динамическую программу только после
предъявления одного кандидата, который проходит общий носитель $P0$, одно
действие и одну нормировку $P1$ и формулирует вычислимый тест запуска $P2$.
Настоящий гейт проверяет эти три пункта для поля
\begin{equation}
 \Phi\in\Gamma\!\left(
 X,
 E_{\rm aff}\widehat\otimes\mathcal Y_{\rm phys}
 \right),
 \qquad
 E_{\rm aff}=\operatorname{Hom}(\mathbb C^4,\operatorname{im}P_3).
\end{equation}

Проверка не ищет массу, частицу или наблюдаемое совпадение. Она отвечает на
более ранний вопрос: является ли кандидат новым общим родителем либо скрытым
повтором закрытых конструкций.

\section{Сверка с закрытыми маршрутами}

Кандидат отличается от четырёх уже проверенных схем.
\begin{enumerate}
 \item В отличие от обменного моста $K$, поле $\Phi$ не является одной
 заранее выбранной матричной единицей между двумя ориентациями. Оно пробегает
 полный физически типизированный модуль одноформ.
 \item В отличие от проекторной модели $Q$, семейное состояние
 $R_\Phi$ является производной квадратичной характеристикой $\Phi$, а не
 независимым входом действия.
 \item В отличие от факторизованного подъёма
 $I_3\otimes U_{\rm walk}$, поле не действует тождественно в семейном
 факторе: его аффинная компонента сама является динамической линейной картой.
 \item В отличие от compacton-ветви, локальное правило обновления и его фаза
 не задаются заранее. Пространственная динамика должна следовать из той же
 суперсвязности и нормы кривизны.
\end{enumerate}

Наиболее близким ранним объектом является градуированное соответствие Тома V.
Новое содержание состоит в том, что аффинный ранговый модуль, физическая
одноформа и пространственная связность объединяются до выбора проекторного
состояния. Суперсвязности как общий математический язык стандартны
\cite{v7-quillen1985}; проектная новизна относится не к их существованию, а к
конкретному входному контракту и запрету обратного проектирования.

\section{Проверка общего носителя \texorpdfstring{$P0$}{P0}}

Пусть
\begin{equation}
 P_1=\frac14J_4,
 \qquad
 P_3=I_4-P_1.
\end{equation}
Тогда $\rank P_3=3$ и
\begin{equation}
 \dim_{\mathbb C}E_{\rm aff}=3\cdot4=12.
\end{equation}
Поле имеет двадцать четыре вещественных аффинных компонент до умножения на
физический модуль одноформ. Действие перестановки $g\in S_4$ имеет вид
\begin{equation}
 \Phi\longmapsto
 \rho_3(g)\Phi\rho_4(g)^{-1},
\end{equation}
поэтому пространство полей инвариантно, хотя отдельная конфигурация может
динамически нарушать $S_4$.

Хиральная и калибровочная типизация берутся не из произвольного полного
эндоморфизма, а из уже выделенного модуля
$\mathcal Y_{\rm phys}=\Omega^1_{D_F}(\mathcal A_F)$. Пространственная
зависимость задаётся сечением над $X$, а Real-завершение --- суммой
$\Phi+\varepsilon'J\Phi J^{-1}$. Следовательно, четыре требуемых фактора
присутствуют в одном поле, а не в прямой сумме четырёх независимых действий.

\begin{proposition}[Предварительный проход $P0$]
Поле $\Phi$ задаёт общий типизированный носитель семейного, хирального,
пространственного и Real-секторов. Его ненулевые элементы являются
нескалярными интертвинерами по аффинному фактору. Полный физический размер и
junk-фактор модуля $\mathcal Y_{\rm phys}$ должны быть повторно вычислены в
следующем гейте, но прямое факторизованное тождество не является определением
кандидата.
\end{proposition}

\section{Проверка одного действия \texorpdfstring{$P1$}{P1}}

Рассмотрим
\begin{equation}
 \mathbb A_\Phi=\mathbb A_0+\widehat\Phi,
 \qquad
 \mathbb F_\Phi=\mathbb A_\Phi^2
\end{equation}
и один функционал
\begin{equation}
 \mathcal S_7
 =\int_X\operatorname{tr}_{\rm norm}
 (\mathbb F_\Phi^*\mathbb F_\Phi)d\mu_X.
 \label{eq:v7-admission-single-action}
\end{equation}
Нормированный след на тензорном произведении определяется произведением
нормированных конечных следов и пространственного $L^2$-скалярного
произведения. Отдельные веса для семейной, Real-, хиральной и исходящей мод
не разрешены.

Нулевой сектор, ранговые страты, локальные поля и линейные исходящие
возмущения являются конфигурациями или касательными одного функционала.
Это ещё не доказывает правильность физической меры, но устраняет прежнюю
склейку нескольких несовместимых действий.

\begin{proposition}[Предварительный проход $P1$]
Формула \eqref{eq:v7-admission-single-action} задаёт один родительский
функционал с одним следом и одной системой размерностей. На входном уровне
$P1$ пройден как спецификация кандидата. Любое последующее появление
независимого секторного коэффициента закрывает этот проход.
\end{proposition}

\section{Физический тест запуска \texorpdfstring{$P2$}{P2}}

Разложение кривизны имеет вид
\begin{equation}
 \mathbb F_{t\eta}
 =\mathbb F_0+tL_0\eta+t^2\widehat\eta^{,2},
 \qquad
 L_0\eta=[\mathbb A_0,\widehat\eta]_{\rm s}.
\end{equation}
После проверки стационарности нулевого сектора физический гессиан равен
\begin{equation}
 H_0(\eta)
 =2\|L_0\eta\|^2
 +4\operatorname{Re}
 \langle\mathbb F_0,\widehat\eta^{,2}\rangle.
 \label{eq:v7-admission-hessian}
\end{equation}
Тем самым $P2$ является конечной спектральной задачей после выбора уже
зафиксированных представлений и факторирования нефизических направлений.

\section{Минимальный конечномерный контроль}

Для проверки алгебраической части берётся ортонормированный базис
sum-zero подпространства $\mathbb R^4$ и прямоугольное поле
$X\in M_{3\times4}(\mathbb C)$. Соответствующий нечётный оператор равен
\begin{equation}
 B(X)=
 \begin{pmatrix}
  0&X^*\\
  X&0
 \end{pmatrix}
 \quad\text{на }\mathbb C^4\oplus\mathbb C^3.
\end{equation}
Он самосопряжён, нечётен относительно
$\Gamma=\operatorname{diag}(I_4,-I_3)$ и имеет
\begin{equation}
 \rank B(X)=2\rank X.
\end{equation}
Real-удвоение $B(X)\oplus\overline{B(X)}$ удовлетворяет обменному
Real-условию точно.

В плоском минимальном фоне $\mathbb A_0=\Gamma$ имеем
\begin{equation}
 \{\Gamma,B(X)\}=0,
 \qquad
 \mathbb F_0=I_7.
\end{equation}
Поэтому на всех двадцати четырёх вещественных направлениях
\begin{equation}
 \Hess_0\mathcal S_7=\frac87I_{24}.
 \label{eq:v7-flat-surrogate-hessian}
\end{equation}

\begin{theorem}[Нулевая граница минимального фона]
Плоский аффинно-хиральный суррогат кандидата не имеет отрицательной моды.
Следовательно, сам факт переменного ранга не запускает переход. Если полный
родитель Тома VII неустойчив, отрицательный вклад обязан происходить из
нецентральной кривизны физического $D_F$, пространственной связи или их
канонического смешанного члена в \eqref{eq:v7-admission-hessian}.
\end{theorem}

Это отрицательный контроль, а не провал всего кандидата: входное требование
$P2$ состояло в формулировке вычислимого теста. Но оно запрещает вводить после
этого ручную отрицательную массу для спасения перехода.

\section{Решение о допуске}

\begin{center}
\begin{tabular}{p{0.12\textwidth}p{0.19\textwidth}p{0.57\textwidth}}
\toprule
Пункт & Статус & Основание \\
\midrule
$P0$ & предварительно пройден & одно поле содержит аффинный, физически
хиральный, пространственный и Real-факторы \\
$P1$ & предварительно пройден & одна суперсвязность, одна норма кривизны,
один нормированный след \\
$P2$ & сформулирован & физический гессиан задан
\eqref{eq:v7-admission-hessian}; плоский контроль строго положителен \\
$P3$--$P5$ & не проверены & конечный объект, скорость и масштаб запрещено
обсуждать до полного $P2$ \\
$P6$ & пройден для входа & гейт, аудит, результат, нулевая гипотеза и
стоп-критерий созданы одновременно \\
\bottomrule
\end{tabular}
\end{center}

\begin{theorem}[Допуск Тома VII]
Рангоизменяющая суперсвязность проходит входной контракт открытия нового
тома на уровне $P0$, $P1$ и вычислимости $P2$. Это допускает исследование
полного физического гессиана, но не является доказательством рождения
материи. Минимальный плоский фон стабилен, поэтому следующий гейт обязан
использовать полный ранее зафиксированный физический модуль без добавленной
отрицательной массы.
\end{theorem}

\begin{nogocriterion}[Стоп-критерий родителя Тома VII]
Кандидат закрывается, если после полного gauge-, Real-, junk- и
BRST/BV-факторирования пространство допустимых $\Phi$ пусто или факторизовано,
единый след требует независимых относительных весов либо оператор
\eqref{eq:v7-admission-hessian} неотрицателен на полном физическом носителе.
В этом случае запрещено переходить к compacton, нити, протону или формулам
массы.
\end{nogocriterion}

\begin{protocol}
\begin{enumerate}
 \item аффинный проектор: \textbf{$\rank P_3=3$};
 \item комплексная размерность $E_{\rm aff}$: \textbf{$12$};
 \item вещественных аффинных вариаций: \textbf{$24$};
 \item ранги $B(X)$ для $\rank X=0,1,2,3$: \textbf{$0,2,4,6$};
 \item ковариантность относительно $S_4$: \textbf{проверена};
 \item самосопряжённость и Real-условие: \textbf{проверены};
 \item один функционал: \textbf{задан};
 \item минимальный гессиан: \textbf{$(8/7)I_{24}$};
 \item отрицательных мод плоского суррогата: \textbf{ноль};
 \item следующий гейт: \textbf{полный физический модуль и гессиан
 нецентральной кривизны}.
\end{enumerate}
\end{protocol}

Машинный сертификат сохраняется в
\path{s2t_v7_rank_changing_superconnection_admission_gate_results.json}.

\begin{thebibliography}{9}
\bibitem{v7-quillen1985}
D.~Quillen,
\emph{Superconnections and the Chern Character},
Topology 24 (1985), 89--95.
\end{thebibliography}